\documentclass{article}
\usepackage{amsmath, amssymb, amsthm}
\usepackage{algorithm}
\usepackage{algpseudocode}
\usepackage{geometry}
\geometry{a4paper, margin=1in}

\newtheorem{theorem}{Theorem}
\newtheorem{lemma}[theorem]{Lemma}
\newtheorem{definition}[theorem]{Definition}

\begin{document}

\section*{Randomized Coordinate Descent (RCD)}
\section*{Mathematical Formulation}
Let $f: \mathbb{R}^d \to \mathbb{R}$ be a differentiable non-convex objective function. At iteration $t$, the algorithm selects a coordinate $i_t \in \{1, 2, \dots, d\}$ randomly with probability $p_{i_t} > 0$, where $\sum_{i=1}^d p_i = 1$. Let $e_i$ denote the $i$-th standard basis vector in $\mathbb{R}^d$. The update rule is given by:
\begin{align*}
w_{t+1} &= w_t - \eta \frac{1}{p_{i_t}} \nabla_{i_t} f(w_t) e_{i_t}
\end{align*}
where $\eta > 0$ is the learning rate (step size), and $\nabla_{i} f(w)$ denotes the partial derivative of $f$ with respect to the $i$-th coordinate.

\section*{Pseudocode}
\begin{algorithm}[H]
\caption{Randomized Coordinate Descent for Non-Convex Optimization}
\begin{algorithmic}[1]
\State \textbf{Input:} Initial point $w_0 \in \mathbb{R}^d$, step size $\eta > 0$, iteration count $T$, probability distribution $\{p_i\}_{i=1}^d$
\For{$t = 0, 1, \dots, T-1$}
    \State Sample coordinate $i_t \in \{1, \dots, d\}$ with probability $\mathbb{P}(i_t = i) = p_i$
    \State Compute partial derivative $g_{i_t} = \nabla_{i_t} f(w_t)$
    \State Update $w_{t+1} = w_t - \eta \frac{1}{p_{i_t}} g_{i_t} e_{i_t}$
\EndFor
\State \textbf{Return:} $\hat{w}$ sampled uniformly at random from $\{w_0, w_1, \dots, w_{T-1}\}$
\end{algorithmic}
\end{algorithm}

\section*{Dry Run Example}
Consider the objective function $f(w) = (w-3)^2$ with initial value $w_0 = 0$, step size $\eta = 0.1$, and $d=1$. Since there is only one coordinate, $i_t = 1$ with probability $p_1 = 1$.
\begin{itemize}
    \item \textbf{Iteration 1:}
    \begin{align*}
        g_1 &= \nabla_1 f(w_0) = 2(w_0 - 3) = 2(0 - 3) = -6 \\
        w_1 &= w_0 - \eta \frac{1}{p_1} g_1 e_1 = 0 - 0.1 \times \frac{1}{1} \times (-6) \times 1 = 0.6 \\
        f(w_1) &= (0.6 - 3)^2 = (-2.4)^2 = 5.76
    \end{align*}
    \item \textbf{Iteration 2:}
    \begin{align*}
        g_1 &= \nabla_1 f(w_1) = 2(w_1 - 3) = 2(0.6 - 3) = -4.8 \\
        w_2 &= w_1 - \eta \frac{1}{p_1} g_1 e_1 = 0.6 - 0.1 \times \frac{1}{1} \times (-4.8) \times 1 = 1.08 \\
        f(w_2) &= (1.08 - 3)^2 = (-1.92)^2 = 3.6864
    \end{align*}
    \item \textbf{Iteration 3:}
    \begin{align*}
        g_1 &= \nabla_1 f(w_2) = 2(w_2 - 3) = 2(1.08 - 3) = -3.84 \\
        w_3 &= w_2 - \eta \frac{1}{p_1} g_1 e_1 = 1.08 - 0.1 \times \frac{1}{1} \times (-3.84) \times 1 = 1.464 \\
        f(w_3) &= (1.464 - 3)^2 = (-1.536)^2 = 2.3593
    \end{align*}
\end{itemize}

\section*{Assumptions}
\begin{definition}[Coordinate-wise Smoothness]
The function $f$ is $L$-coordinate-wise smooth if for all $w \in \mathbb{R}^d$, all coordinates $i \in \{1, \dots, d\}$, and any $\alpha \in \mathbb{R}$, the following inequality holds:
\begin{align*}
    f(w + \alpha e_i) \leq f(w) + \alpha \nabla_i f(w) + \frac{L}{2} \alpha^2
\end{align*}
\end{definition}
\begin{definition}[Bounded Variance]
The stochastic gradient variance is bounded such that there exists a constant $\sigma \geq 0$ where:
\begin{align*}
    \mathbb{E}_{i} \left[ \left\| \frac{1}{p_i} \nabla_i f(w) e_i - \nabla f(w) \right\|^2 \right] \leq \sigma^2
\end{align*}
\end{definition}
\textbf{Remark:} For uniform sampling $p_i = 1/d$, $\sigma^2 = (d-1)\|\nabla f(w)\|^2$. For importance sampling $p_i \propto |\nabla_i f(w)|$, $\sigma^2 = 0$.

\section*{Main Convergence Theorem}
\begin{theorem}
Let $f$ be $L$-coordinate-wise smooth and lower bounded by $f^*$. Let $\eta \leq \frac{1}{L \max_i \frac{1}{p_i}}$. After $T$ iterations of RCD, the algorithm satisfies:
\begin{align*}
    \frac{1}{T} \sum_{t=0}^{T-1} \mathbb{E}[\|\nabla f(w_t)\|^2] \leq \frac{2(f(w_0) - f^*)}{\eta T} + \frac{L \eta \sigma^2}{2}
\end{align*}
By setting $\eta = \min\left( \frac{1}{L \max_i \frac{1}{p_i}}, \frac{2(f(w_0)-f^*)}{L \sigma^2 T} \right)$, the convergence rate is:
\begin{align*}
    \frac{1}{T} \sum_{t=0}^{T-1} \mathbb{E}[\|\nabla f(w_t)\|^2] \leq O\left(\frac{1}{T} + \frac{\sigma}{\sqrt{T}}\right)
\end{align*}
\end{theorem}

\section*{Theoretical Properties}
\begin{itemize}
    \item \textbf{Stability:} The step size $\eta$ is constrained by the coordinate-wise smoothness constant and the maximum inverse probability $1/p_i$. This ensures that the expected descent dominates the variance-induced second-order error.
    \item \textbf{Hyperparameter Sensitivity:} If $p_i$ is set too small for a coordinate with large gradient, $1/p_i$ explodes, severely restricting the allowable step size $\eta$ and slowing convergence. Importance sampling (choosing $p_i \propto |\nabla_i f(w)|$) eliminates the variance term ($\sigma=0$) and allows larger step sizes, recovering deterministic rates.
    \item \textbf{Comparison to Baselines:} Compared to full-batch Gradient Descent (GD), RCD uses only $1/d$ of the gradient computation per step but suffers from variance. When $\sigma > 0$, RCD converges at $O(1/\sqrt{T})$ to a neighborhood, whereas GD converges at $O(1/T)$. However, per-iteration cost is $d$ times cheaper.
\end{itemize}

\section*{Discussion}
The $1/p_i$ weighting in the update rule is strictly necessary to ensure that the expected update equals the full gradient: $\mathbb{E}[\frac{1}{p_{i_t}} \nabla_{i_t} f(w_t) e_{i_t}] = \nabla f(w_t)$. Without this correction, the algorithm would exhibit biased convergence, stalling at non-stationary points. The variance bound $\sigma^2$ scales with $\max_i 1/p_i$, underlining the critical role of the sampling distribution.

\section*{Preliminaries (Definitions and Lemmas)}
\begin{lemma}[Descent Lemma for Coordinate Update]
If $f$ is $L$-coordinate-wise smooth, then for any $w \in \mathbb{R}^d$, coordinate $i$, and step size $\alpha$:
\begin{align*}
    f(w - \alpha \nabla_i f(w) e_i) \leq f(w) - \alpha \|\nabla_i f(w)\|^2 + \frac{L}{2} \alpha^2 \|\nabla_i f(w)\|^2
\end{align*}
\end{lemma}
\begin{proof}
Apply the coordinate-wise smoothness definition with $\alpha \nabla_i f(w)$ as the scalar displacement:
\begin{align*}
    f(w - \alpha \nabla_i f(w) e_i) &\leq f(w) + (-\alpha \nabla_i f(w)) \nabla_i f(w) + \frac{L}{2} (-\alpha \nabla_i f(w))^2 \\
    &= f(w) - \alpha \|\nabla_i f(w)\|^2 + \frac{L}{2} \alpha^2 \|\nabla_i f(w)\|^2
\end{align*}
\end{proof}

\begin{lemma}[Variance Decomposition]
For any $w \in \mathbb{R}^d$, the second moment of the scaled coordinate gradient equals the squared norm of the full gradient plus the variance:
\begin{align*}
    \mathbb{E}_{i} \left[ \left\| \frac{1}{p_i} \nabla_i f(w) e_i \right\|^2 \right] = \|\nabla f(w)\|^2 + \sigma^2
\end{align*}
\end{lemma}
\begin{proof}
By definition of variance:
\begin{align*}
    \mathbb{E}_{i} \left[ \left\| \frac{1}{p_i} \nabla_i f(w) e_i - \nabla f(w) \right\|^2 \right] &= \mathbb{E}_{i} \left[ \left\| \frac{1}{p_i} \nabla_i f(w) e_i \right\|^2 \right] - \|\nabla f(w)\|^2
\end{align*}
Rearranging yields the result.
\end{proof}

\section*{Main Proof (Step-by-Step Derivation)}
We derive the expected decrease in the objective function at iteration $t$.
\begin{align*}
    &\mathbb{E}_{i_t}[f(w_{t+1})] \\
    &= \mathbb{E}_{i_t}[f(w_t - \eta \frac{1}{p_{i_t}} \nabla_{i_t} f(w_t) e_{i_t})] \tag{Substitute update rule} \\
    &\leq \mathbb{E}_{i_t} \left[ f(w_t) - \eta \frac{1}{p_{i_t}} \nabla_{i_t} f(w_t) \cdot \nabla_{i_t} f(w_t) + \frac{L}{2} \eta^2 \frac{1}{p_{i_t}^2} (\nabla_{i_t} f(w_t))^2 \right] \tag{Apply Lemma 1 with $\alpha = \eta / p_{i_t}$} \\
    &= f(w_t) - \eta \mathbb{E}_{i_t} \left[ \frac{1}{p_{i_t}} \|\nabla_{i_t} f(w_t)\|^2 \right] + \frac{L \eta^2}{2} \mathbb{E}_{i_t} \left[ \frac{1}{p_{i_t}^2} \|\nabla_{i_t} f(w_t)\|^2 \right] \tag{Linearity of Expectation} \\
    &= f(w_t) - \eta \sum_{i=1}^d p_i \frac{1}{p_i} \|\nabla_i f(w_t)\|^2 + \frac{L \eta^2}{2} \sum_{i=1}^d p_i \frac{1}{p_i^2} \|\nabla_i f(w_t)\|^2 \tag{Expand expectation} \\
    &= f(w_t) - \eta \sum_{i=1}^d \|\nabla_i f(w_t)\|^2 + \frac{L \eta^2}{2} \sum_{i=1}^d \frac{1}{p_i} \|\nabla_i f(w_t)\|^2 \tag{Cancel $p_i$ terms} \\
    &= f(w_t) - \eta \|\nabla f(w_t)\|^2 + \frac{L \eta^2}{2} \mathbb{E}_{i_t} \left[ \left\| \frac{1}{p_{i_t}} \nabla_{i_t} f(w_t) e_{i_t} \right\|^2 \right] \tag{Sum is full gradient norm; re-write as expectation} \\
    &= f(w_t) - \eta \|\nabla f(w_t)\|^2 + \frac{L \eta^2}{2} \left( \|\nabla f(w_t)\|^2 + \sigma^2 \right) \tag{Apply Lemma 2} \\
    &= f(w_t) - \left( \eta - \frac{L \eta^2}{2} \right) \|\nabla f(w_t)\|^2 + \frac{L \eta^2}{2} \sigma^2 \tag{Group gradient terms} \\
    &\leq f(w_t) - \frac{\eta}{2} \|\nabla f(w_t)\|^2 + \frac{L \eta^2}{2} \sigma^2 \tag{Assume $\eta \leq 1/L$, so $\eta - L\eta^2/2 \geq \eta/2$}
\end{align*}
Taking total expectation over all randomness $\mathcal{F}_t = \{i_0, \dots, i_{t-1}\}$:
\begin{align*}
    \mathbb{E}[f(w_{t+1})] &\leq \mathbb{E}[f(w_t)] - \frac{\eta}{2} \mathbb{E}[\|\nabla f(w_t)\|^2] + \frac{L \eta^2}{2} \sigma^2 \tag{Total expectation}
\end{align*}

\section*{Rate Derivation (With Constants)}
Rearrange the total expectation inequality to isolate the gradient norm term:
\begin{align*}
    \frac{\eta}{2} \mathbb{E}[\|\nabla f(w_t)\|^2] &\leq \mathbb{E}[f(w_t)] - \mathbb{E}[f(w_{t+1})] + \frac{L \eta^2}{2} \sigma^2 \tag{Rearrange}
\end{align*}
Sum this inequality from $t=0$ to $T-1$:
\begin{align*}
    \frac{\eta}{2} \sum_{t=0}^{T-1} \mathbb{E}[\|\nabla f(w_t)\|^2] &\leq \mathbb{E}[f(w_0)] - \mathbb{E}[f(w_T)] + \frac{T L \eta^2}{2} \sigma^2 \tag{Telescoping sum} \\
    &\leq f(w_0) - f^* + \frac{T L \eta^2}{2} \sigma^2 \tag{Lower bound $f(w_T) \geq f^*$}
\end{align*}
Divide by $T \eta / 2$:
\begin{align*}
    \frac{1}{T} \sum_{t=0}^{T-1} \mathbb{E}[\|\nabla f(w_t)\|^2] &\leq \frac{2(f(w_0) - f^*)}{\eta T} + L \eta \sigma^2 \tag{Divide by $T \eta / 2$}
\end{align*}
To minimize the right hand side with respect to $\eta$, we balance the two terms: $\frac{2(f(w_0)-f^*)}{\eta T} \approx L \eta \sigma^2$.
Solving for $\eta$ yields $\eta^* = \sqrt{\frac{2(f(w_0)-f^*)}{L \sigma^2 T}}$.
Since we also require $\eta \leq \frac{1}{L \max_i 1/p_i}$ to guarantee descent, we set:
\begin{align*}
    \eta = \min \left( \frac{1}{L \max_i \frac{1}{p_i}}, \sqrt{\frac{2(f(w_0)-f^*)}{L \sigma^2 T}} \right)
\end{align*}
\textbf{Case 1:} If $\eta = \frac{1}{L \max_i 1/p_i}$, the bound becomes:
\begin{align*}
    \frac{1}{T} \sum_{t=0}^{T-1} \mathbb{E}[\|\nabla f(w_t)\|^2] \leq \frac{2L \max_i 1/p_i (f(w_0)-f^*)}{T} + \frac{\sigma^2}{\max_i 1/p_i} = O(1)
\end{align*}
This represents convergence to a neighborhood whose radius depends on the variance.
\textbf{Case 2:} If $\eta = \sqrt{\frac{2(f(w_0)-f^*)}{L \sigma^2 T}}$, substituting gives:
\begin{align*}
    \frac{1}{T} \sum_{t=0}^{T-1} \mathbb{E}[\|\nabla f(w_t)\|^2] &\leq \frac{2(f(w_0)-f^*)}{T \sqrt{\frac{2(f(w_0)-f^*)}{L \sigma^2 T}}} + L \sigma^2 \sqrt{\frac{2(f(w_0)-f^*)}{L \sigma^2 T}} \\
    &= \frac{\sqrt{2(f(w_0)-f^*) L \sigma^2}}{\sqrt{T}} + \frac{\sqrt{2(f(w_0)-f^*) L \sigma^2}}{\sqrt{T}} \\
    &= \frac{2\sqrt{2 L \sigma^2 (f(w_0)-f^*)}}{\sqrt{T}} = O\left(\frac{1}{\sqrt{T}}\right)
\end{align*}

\section*{Special Cases}
\textbf{Uniform Sampling:} If $p_i = 1/d$ for all $i$, then $\max_i 1/p_i = d$. The variance is strictly $\sigma^2 = (d-1)\|\nabla f(w)\|^2$ (which is state-dependent). Assuming an absolute bound on gradient norms $\|\nabla f(w)\|^2 \leq G$, we get $\sigma^2 \leq (d-1)G$. The step size constraint becomes $\eta \leq \frac{1}{Ld}$, and the rate degrades by a factor of $d$ compared to standard gradient descent.

\textbf{Importance Sampling:} If we choose $p_i = \frac{|\nabla_i f(w)|}{\|\nabla f(w)\|_1}$, then $\frac{1}{p_i} \nabla_i f(w) = \|\nabla f(w)\|_1 \text{sign}(\nabla_i f(w))$.
The variance evaluates to:
\begin{align*}
    \mathbb{E}_i \left[ \left\| \frac{1}{p_i} \nabla_i f(w) e_i \right\|^2 \right] - \|\nabla f(w)\|^2 = \|\nabla f(w)\|_1^2 - \|\nabla f(w)\|_2^2 \geq 0
\end{align*}
If we use the $L_1$-norm variant of coordinate smoothness, the variance is tightly bounded and often yields better dependency on dimension $d$.

\end{document}